\documentclass[UTF8,fontset=fandol,11pt,a4paper]{ctexart}
\usepackage[a4paper,margin=1.75cm,headheight=15pt]{geometry}
\usepackage{amsmath,amssymb,mathtools}
\usepackage{booktabs,longtable,tabularx,array,multirow}
\usepackage{enumitem}
\usepackage[most]{tcolorbox}
\usepackage{tikz}
\usetikzlibrary{arrows.meta,positioning,fit,calc,matrix}
\usepackage{xcolor}
\usepackage{fancyhdr}
\usepackage{hyperref}
\usepackage{microtype}
\usepackage{lastpage}

\newcolumntype{Y}{>{\raggedright\arraybackslash}X}
\newcolumntype{L}[1]{>{\raggedright\arraybackslash}p{#1}}
\newcolumntype{C}[1]{>{\centering\arraybackslash}p{#1}}
\setlength{\parindent}{2em}
\setlength{\parskip}{0.34em}
\setlength{\emergencystretch}{3em}
\renewcommand{\arraystretch}{1.2}
\setlength{\tabcolsep}{3pt}
\setlist[itemize]{itemsep=0.18em,topsep=0.35em,leftmargin=2em}
\setlist[enumerate]{itemsep=0.18em,topsep=0.35em,leftmargin=2em}

\definecolor{deep}{RGB}{38,58,72}
\definecolor{soft}{RGB}{244,247,248}
\definecolor{accent}{RGB}{116,72,46}
\definecolor{greenish}{RGB}{52,91,74}
\definecolor{warnc}{RGB}{136,61,58}
\definecolor{purple}{RGB}{84,72,112}
\definecolor{rulegray}{RGB}{110,116,120}

\hypersetup{colorlinks=true,linkcolor=deep,urlcolor=accent,pdftitle={CO-02 ISA 与机器级程序}}
\pagestyle{fancy}
\fancyhf{}
\lhead{\small\color{deep}\textbf{408 · 计算机组成原理 CO-02}}
\rhead{\small\color{deep}ISA · 机器语义 · ABI}
\cfoot{\small 第 \thepage\ 页 / 共 \pageref{LastPage} 页}
\renewcommand{\headrulewidth}{0.4pt}

\newtcolorbox{corebox}[1]{breakable,colback=soft,colframe=deep,title=\textbf{#1},boxrule=0.8pt,arc=2mm}
\newtcolorbox{methodbox}[1]{breakable,colback=white,colframe=greenish,title=\textbf{#1},boxrule=0.7pt,arc=1.5mm}
\newtcolorbox{warnbox}[1]{breakable,colback=white,colframe=warnc,title=\textbf{#1},boxrule=0.7pt,arc=1.5mm}
\newtcolorbox{examplebox}[1]{breakable,colback=white,colframe=purple,title=\textbf{#1},boxrule=0.7pt,arc=1.5mm}
\newcommand{\kw}[1]{\textbf{\color{deep}#1}}
\newcommand{\code}[1]{\texttt{#1}}

\tikzset{
  co/.style={draw=deep,rounded corners=1.5mm,text width=2.8cm,minimum height=0.95cm,align=center,fill=soft,font=\small,inner sep=3pt},
  state/.style={co,double,double distance=1.2pt,fill=white},
  abstract/.style={co,fill=black!8},
  flow/.style={-{Latex[length=2mm]},thick,draw=deep},
  ref/.style={-{Latex[length=2mm]},dashed,draw=rulegray}
}

\title{\vspace{-1.1cm}\textbf{ISA 与机器级程序}\\[0.4em]
\large 从程序意图到下一架构状态的可执行契约}
\author{408 · 计算机组成原理 Topic CO-02}
\date{}

\begin{document}
\maketitle

\begin{corebox}{母问题：软件意图怎样成为处理器必须兑现的语义？}
高级语言只说值、控制流和存储访问想要什么；处理器必须把它们编码成有限长度的指令，并规定执行后软件能观察到的下一状态。本册的生成链是
\[
\boxed{\text{Program Intent}
\to\text{Architectural State}
\to\text{Encoding / Address Semantics}
\to\text{Next Architectural State}}
\]
箭头首先表示语义约束的生成方向：意图决定所需状态变化，状态变化必须有编码和地址解释，最后形成新的架构状态。它不是某张具体数据通路图的时序；CO-03 才负责把已经确定的语义交给通路与控制器实现。
\end{corebox}

\begin{methodbox}{本册学习范围}
\begin{itemize}
  \item 本册说明指令集如何规定寄存器、程序计数器和内存的可观察变化，并把这些变化编码成机器指令。
  \item 本册解释操作码、寄存器字段、立即数、有效地址、访问宽度、字节序、对齐和函数调用约定；遇到具体处理器时，题设优先。
  \item 流水线、Cache 和页表只在需要说明指令语义的地方定义；读者不必先读其他册才能理解本册的基本推导。
\end{itemize}
\end{methodbox}

\section{术语先行：一条机器指令究竟承诺什么}
\begin{methodbox}{最常用的八个词}
\begin{itemize}
  \item \kw{ISA（指令集体系结构）：}软件可观察的指令、寄存器、地址和异常规则；它描述“必须发生什么”，不规定内部电路怎样连接。
  \item \kw{架构状态：}执行指令前后软件能够观察到的 PC、寄存器、内存和状态寄存器内容。
  \item \kw{PC（程序计数器）：}保存下一条取指位置的寄存器；分支、跳转、调用和异常可以把它改为非顺序地址。
  \item \kw{有效地址 EA：}真正用于一次访存的地址，例如基址寄存器加位移；它是地址，不是地址处存放的数据。
  \item \kw{ABI（应用二进制接口）：}编译器、函数和目标文件之间约定参数、返回值、寄存器保存责任和栈布局的规则。
  \item \kw{字节序：}一个多字节对象分布到连续地址时，低位字节和高位字节的排列顺序。
  \item \kw{对齐：}对象起始地址满足特定位数边界，例如 4 字节对象从 4 的倍数地址开始。
  \item \kw{RISC/CISC：}两组常见的指令集设计取舍；它们不是“硬布线/微程序”的必然对应关系。
\end{itemize}
\end{methodbox}

\begin{center}
\begin{tikzpicture}[node distance=8mm and 5mm]
  \node[abstract] (intent) {程序意图\\值、控制流、访问};
  \node[state,right=of intent] (arch) {架构状态\\PC、寄存器、内存};
  \node[co,right=of arch] (enc) {编码与地址语义\\字段、EA、宽度};
  \node[state,right=of enc] (next) {下一状态\\结果、PC、异常};
  \draw[flow] (intent)--(arch);
  \draw[flow] (arch)--(enc);
  \draw[flow] (enc)--(next);
  \draw[ref] (next.south) to[bend right=24] node[below,font=\scriptsize]{语义等价的下一条指令} (intent.south);
\end{tikzpicture}
\end{center}

\tableofcontents
\newpage

\section{先固定对象：架构状态不是硬件内部状态}

\subsection{软件契约层}
设架构状态为
\[
A=(PC,R,M,F,C,\ldots),
\]
其中 $PC$ 是程序计数器，$R$ 是架构寄存器，$M$ 是软件可观察的内存效果，$F$ 是架构标志，$C$ 表示 ISA 规定的控制/特权状态。指令 $I$ 定义状态转移
\[
A_{t+1}=\mathcal I(A_t).
\]
ISA 的稳定问题是：\kw{执行后哪些状态必须改变，哪些状态必须保持不变，以及异常发生时允许观察到什么。}

\subsection{实现层}
IR、MAR、MDR、临时寄存器、流水寄存器和 Cache 元数据可以是微架构状态。它们服务于实现，却不自动成为 ISA 规定的对象。相同 ISA 可以使用不同的寄存器数量、流水级数、内部总线和控制器，只要架构状态轨迹与契约等价。

\begin{warnbox}{不可逆推的边界}
看到教材图上的固定 PC 加法器、MAR/MDR 或某个九拍节奏，不能反推出所有 ISA 都必须有这些部件。ISA 定义“必须发生什么”，CO-03 才讨论“一种实现怎样让它发生”。
\end{warnbox}

\subsection{可访问性是 ISA 相对的}
“用户可见/部分可见/系统可见/完全不可见”只能作为提问方向，不是跨架构固定清单。稳定判据是：当前 ISA 是否提供读取、写入、配置或间接影响该状态的合法操作？例如 PC 可能被跳转间接改变，页表基址可能有特权接口，PCB 通常是操作系统软件对象而非必然 CPU 寄存器；Cache 对普通指令透明，却可能有特权管理操作。

\section{指令编码是有限预算分配}

\subsection{四个最小问题}
每条指令必须回答：\kw{做什么}（opcode）、\kw{对谁做}（操作数/地址）、\kw{结果放哪}（目的状态）、\kw{下一步去哪}（顺序 PC 或控制转移目标）。地址字段数量形成零地址、一地址、二地址、三地址等表示，但地址数量不是性能或表达能力的单一排名。

\subsection{扩展操作码的生成}
若指令字长为 $L$，本层保留 $w$ 位地址字段，操作码拥有 $L-w$ 位。已经使用的编码不能再分配，剩余前缀才可以作为下一层的 escape：
\[
N_{k+1}=\bigl(N_k-U_k\bigr)2^{w_k},
\]
其中 $N_k$ 是本层可用编码数，$U_k$ 是本层已经占用的指令数，$w_k$ 是被释放给下一层的字段位数。

\begin{examplebox}{16 位字、6 位地址字段}
二地址层保留 $4$ 位操作码，共 $16$ 个编码。若已用 $15$ 个，剩下 $1$ 个作为扩展前缀；一地址层得到 $1\times2^6=64$ 个候选编码，若已用 $34$ 个，零地址层只从剩余 $30\times2^6=1920$ 个编码中继续展开。每一层都必须同时检查“剩余数量”和“前缀无歧义”，不能把 escape 当普通指令再算一次。
\end{examplebox}

\subsection{定长与变长、RISC 与 CISC}
定长格式让字段位置和取指边界规整，便于译码和并行取指；变长格式可以提高代码密度，却增加边界识别和译码成本。RISC/CISC 是一组设计权衡：指令规整度、内存操作约束、代码密度、译码复杂度、编译器与硬件分工都可能参与，不能推出“RISC 必然硬布线”或“CISC 必然微程序”。

\section{从程序意图到可执行映像：四个转换阶段}
“高级语言能运行”不是编译器一步把文字变成 CPU 动作。稳定的生命周期是：
\[
\text{source}\xrightarrow{\text{compiler}}\text{assembly/relocatable object}
\xrightarrow{\text{assembler}}\text{machine object}
\xrightarrow{\text{linker}}\text{executable image}
\xrightarrow{\text{loader}}\text{process address space}.
\]
编译器把表达式、控制流和函数调用映射为 ISA 可表达的指令与数据布局；汇编器把助记符和伪指令解析为机器码、符号表和重定位记录；链接器解析跨目标文件的符号引用，合并代码/数据段并完成或保留重定位；装入器把可执行映像放入进程地址空间，设置入口 PC、栈和必要的动态绑定。每一步都可能改变表示，但必须保持最终可观察语义。

这条链也解释了三个边界：编译器选择不等于 ISA 必需序列；链接地址不等于运行时物理地址；装入完成不等于第一条指令已经提交。题目给出目标文件、符号或重定位项时，先判断它处在哪一个阶段，再决定谁能修改它。

\begin{center}\small
\begin{tabularx}{\linewidth}{L{2.0cm}YY}
\toprule
阶段 & 主要输入/输出 & 不能直接推出的事实\\\midrule
编译 & 源语言语义 $\to$ 汇编/目标代码 & 唯一指令序列或固定寄存器分配\\
汇编 & 助记符/伪指令 $\to$ 机器码、符号、重定位 & 所有地址已经最终确定\\
链接 & 多个目标文件 $\to$ 可执行映像 & 已经位于某个进程的物理内存\\
装入 & 映像 $\to$ 进程地址空间与入口状态 & 已经完成执行或不会缺页\\
\bottomrule
\end{tabularx}\end{center}

\section{寻址：先形成地址，再访问对象}

\subsection{统一生成式}
有效地址（effective address, EA）是地址形成阶段的输出：
\[
EA=f(\text{instruction fields},R,PC,M).
\]
它回答“对象在哪里”，不等于已经读出了对象。若系统启用虚拟地址，EA 还要经过 CO-07 的翻译接口；若访问 Cache，则继续进入 CO-06。

\begin{center}
\small
\begin{tabularx}{\linewidth}{L{2.2cm}C{1.8cm}C{1.2cm}Y Y}
\toprule
方式 & 操作数/地址来源 & 数据访存 & 语义与题面信号 & 常见误判\\
\midrule
立即 & 字段就是值 & 0 & 常量、符号扩展、位数限制 & 把立即数当地址\\
寄存器 & $R_i$ 中是值 & 0 & 算术和寄存器间数据流 & 把寄存器值再算一次 EA\\
直接 & 字段给 EA & 1 & 形式地址直接定位对象 & 忽略地址空间/编址单位\\
寄存器间接 & $EA=R_i$ & 1 & 指针解引用 & 把指针本身当对象\\
基址+偏移 & $EA=R_b+\operatorname{sext}(d)$ & 1 & 栈帧、结构体、重定位 & 忘记偏移扩展/字节编址\\
变址/比例 & $EA=B+R_iK+d$ & 1 & 数组元素和记录数组 & 忘记 $K=\operatorname{sizeof}(T)$\\
PC 相对 & $target=PC_{base}+d$ & 取指路径 & 分支、位置无关代码 & 不确认 PC 基准和位移单位\\
间接 & 先取地址再解引用 & $\ge2$ & 多级指针或指针表 & 把“取地址”与“取数据”混写\\
\bottomrule
\end{tabularx}
\end{center}

\subsection{地址、值、地址的地址}
做题时把三种量分行写：
\[
\text{value}=M[EA],\qquad
\text{pointer}=M[EA],\qquad
\text{next address}=M[M[EA]].
\]
访存次数必须从这些层级推出来，而不是从寻址名称背口诀。地址形成是一次算术关系；访存是对该地址对应存储对象的读取或写入。

\section{访存宽度、字节序与对齐}

\subsection{字节序}
Endian 只规定多字节对象的字节怎样映射到递增内存地址，不改变该对象在寄存器中代表的数学值。设一个 $32$ 位值由四个字节组成，题目必须先画地址表，再判断低地址存放最高有效字节还是最低有效字节；不能按十六进制字符串的视觉顺序猜。

\subsection{访问宽度与扩展}
Load 的宽度决定读几个字节；signed/unsigned load 决定读入寄存器时如何扩展。Store 只写 ISA 规定宽度的低位或对应字节，不因为寄存器更宽就自动写满整字。地址宽度、指针宽度、PC 宽度、指令长度和物理地址线数属于不同预算，题设没有等式时不能强行相等。

\subsection{对齐}
若对象大小为 $k$ 字节，常见对齐条件是地址满足 $a\equiv0\pmod{k}$，但 ISA 可以允许未对齐访问、将其拆为多次事务，或触发异常。答题必须写清“题设 ISA 的允许行为”，不能把某一实现的惩罚当成普遍语义。

\section{从 C 表达式到机器级状态}

\subsection{数组与表达式}
对元素类型大小 $s=\operatorname{sizeof}(T)$，一维数组地址为
\[
\operatorname{addr}(a[i])=base(a)+i\times s.
\]
行优先二维数组为
\[
\operatorname{addr}(a[i][j])=base+(iN+j)\times s.
\]
这条地址生成式连接三个层次：本册负责语义和 EA，Cache 负责局部性与副本命中，地址翻译负责 VA 到 PA；这里先指出接口，不把后两种机制混进指令语义。

\subsection{控制流}
条件结构通常变为比较、条件选择下一 PC 和无条件跳转形成回边。分支目标可以是 PC 相对、寄存器间接或 ISA 规定的其他形式；比较是否产生显式标志、是否与分支合并，由 ISA 决定。

\subsection{函数调用：ISA 与 ABI 的分层}
函数调用至少有四个状态交接：参数/返回值位置、返回地址、保存责任、栈帧布局。ISA 提供跳转、寄存器和访存等原语；ABI 约定哪些寄存器 caller-saved/callee-saved、参数何时溢出到栈、返回值放在哪里。ABI 不是硬件自动保存规则，具体名称也不能跨架构照抄。

\begin{center}
\begin{tikzpicture}[node distance=7mm and 8mm]
  \node[abstract] (caller) {Caller\\准备参数};
  \node[co,right=of caller] (call) {Call\\更新控制流};
  \node[co,right=of call] (frame) {Callee\\建立栈帧};
  \node[co,right=of frame] (body) {函数体\\读写状态};
  \node[co,below=of body] (ret) {Return\\恢复约定状态};
  \draw[flow] (caller)--(call);
  \draw[flow] (call)--(frame);
  \draw[flow] (frame)--(body);
  \draw[flow] (body)--(ret);
  \draw[ref] (ret) |- (caller.south);
\end{tikzpicture}
\end{center}

\section{贯穿母例：\texorpdfstring{$x=a[i]+1$}{x=a[i]+1}}

设 $a$ 是元素大小为 $s$ 的数组，$i$ 在寄存器中，目标 $x$ 可能是寄存器变量或内存对象。先写架构状态差，而不是先背某套汇编：
\[
\begin{aligned}
EA&=R[base(a)]+R[i]\times s,\\
v&=M[EA],\\
R[x]'&=v+1 \quad\text{或}\quad M[addr(x)]'=v+1.
\end{aligned}
\]
随后才选择一条具体 ISA 的指令序列：计算偏移、形成 EA、load、加立即数、必要时 store。编译器的寄存器分配和优化可能改变序列，但必须保持可观察语义一致。

\begin{warnbox}{母例中的典型错误路径}
“$a[i]$ 已经是一个值，所以直接把 $base+i$ 当地址”漏乘了元素大小；“load 结束，所以 $x$ 已经提交”又把中间值与架构写回混淆。正确顺序是：先确认对象和表示，再形成 EA，再区分产生值与写入目标状态。
\end{warnbox}

\section{不变量、边界与代价}

\subsection{必须保持的不变量}
\begin{itemize}
  \item 同一编码在当前 ISA 下语义确定，不能因某张数据通路图不同而改变；
  \item 指令允许改变的 Write Set 之外，其他架构状态保持不变；
  \item 地址形成、访存宽度、扩展和对齐必须与 ISA 规定一致；
  \item 调用者与被调用者对 ABI 保存责任一致；
  \item 异常或 fault 发生时，架构可见效果符合该 ISA 的精确性规则。
\end{itemize}

\subsection{五列概念边界}
\begin{center}
\small
\begin{tabularx}{\linewidth}{L{1.8cm}C{0.35cm}L{1.8cm}Y Y}
\toprule
概念 A & $\neq$ & 概念 B & 真正区别与题面信号 & 混淆后果\\
\midrule
ISA & $\neq$ & 微架构 & ISA 规定软件可观察语义；微架构选择路径、时序、Cache 与控制器 & 用教材通路否定另一合法实现\\
EA & $\neq$ & value & EA 定位对象；value 是读取对象后的内容 & 少算/多算一次访存\\
ISA & $\neq$ & ABI & ISA 是机器契约；ABI 是编译模块间的调用约定 & 把 caller/callee 保存误当硬件自动动作\\
字节序 & $\neq$ & 位权解释 & 字节序描述内存地址顺序；位权描述位串数值解释 & 反转寄存器值或地址表\\
对齐 & $\neq$ & 地址宽度 & 对齐是合法起始地址约束；宽度是可表示地址/数据范围 & 从 32 位指令推出 32 位地址或固定异常\\
PC 相对 & $\neq$ & 绝对地址 & 相对目标由 PC 基准和位移形成；绝对地址直接给出位置 & 分支范围和重定位判断错误\\
\bottomrule
\end{tabularx}
\end{center}

\subsection{设计权衡}
固定格式、更多寄存器、较大立即数、更多 opcode 会竞争同一编码预算；更复杂的寻址可以减少显式指令，却增加译码和地址形成复杂度；Load/Store 限制让数据通路规整，却可能需要更多指令；寄存器传参减少访存，但寄存器不足时仍需栈帧。任何“更快/更省”的结论都要说明比较对象和成本项。

\section{与相邻机制的接口}

\begin{itemize}
  \item \kw{CO-01：}提供位串解释、有限宽度结果和异常/溢出语义；CO-02 只调用它，不重讲补码电路。
  \item \kw{CO-03：}接收已确定的 Read Set、Derived Values、Write Set 和 Next PC，生成通路、调度与控制字。
  \item \kw{CO-04：}接收单指令依赖和语义，负责多指令时间重叠、forwarding、stall、flush 与 ordered commit。
  \item \kw{CO-06/07：}CO-02 只输出访问宽度、EA/VA 和重试语义；Cache 命中状态、TLB/page walk 和页表硬件不在本册拥有。
  \item \kw{OS / X-B01：}异常入口、特权切换和 handler 的软件动作是跨层接口；不能把“用户不可见”简化成“所有特权软件都不可见”。
\end{itemize}

\section{做题控制协议}

看到指令格式、寻址、机器级代码、栈帧或字节序题，按以下顺序执行：
\begin{enumerate}
  \item 写题设 ISA、编址单位、指令长度和操作数宽度；缺失时列出不能强推的量；
  \item 写指令前后的架构状态差：Read Set、Write Set、Next PC、异常可能性；
  \item 解码字段预算，明确 opcode、register、immediate 和 function 各占哪些位；
  \item 对每个内存操作先写 $EA$ 生成式，再区分 value、pointer、address-of-address 和访存次数；
  \item 若是数组/栈帧，先乘元素大小或确认偏移单位；若是多字节对象，画地址—字节表；
  \item 函数题分开 ISA 原语、ABI 约定和编译器选择；不要把某架构寄存器名称当成普适规则；
  \item 最后把语义交给 CO-03/CO-04，独立检查状态差、地址单位、写回位置和异常可见效果。
\end{enumerate}

\begin{methodbox}{压缩成第一动作}
题面出现“指令格式/寻址/机器代码/调用/大小端”时，先问：\kw{这条指令要改变哪个架构状态？操作数现在是值、地址还是地址的地址？} 写出状态差和 EA 后，剩余计算才有唯一落点。
\end{methodbox}

\section{本册与硬件实现的接口}
\begin{corebox}{从架构状态到数据位置}
本册先写指令前后的架构状态差，再由编码、寻址和访问宽度说明操作数如何被定位。输出是一组完整的读集合、写集合、下一条 PC 和异常条件，硬件可以据此设计通路，但不能把某套微架构的内部节拍当成指令集事实。
\end{corebox}
本册局部成本来自编码预算、指令长度、EA 生成、对齐/字节序处理与 ABI 约定；它们可能改变取指、访存和调用开销，但程序总时间仍需结合 IC、CPI、时钟周期及 Cache/Memory 行为。

\section{章节与考纲映射}

\begin{tabularx}{\linewidth}{L{3.2cm}L{5.4cm}Y}
\toprule
传统知识块 & 本册的机制位置 & 本册不展开的实现细节\\
\midrule
指令系统与指令格式 & 编码预算、字段竞争、扩展操作码 & 具体控制器实现归 CO-03\\
寻址方式与有效地址 & EA 生成、地址/值分层、访存次数 & 翻译硬件归 CO-07\\
机器级程序与 C 映射 & 数组地址、控制流、函数调用、栈帧 & 优化器内部选择只作实现例子\\
RISC/CISC & 代码密度、规整译码、内存访问和软硬件分工的权衡 & 不推出控制器一一对应\\
字节序、对齐、访问宽度 & 地址表、扩展、合法访问边界 & 平台特例必须由题设声明\\
\bottomrule
\end{tabularx}

\section{一页压缩与自测}

\begin{corebox}{一句话}
ISA 先规定“程序执行后软件必须看到什么”；机器级程序把意图编码成状态变化，寻址把字段解释为对象位置，ABI 把调用双方接起来，而 CO-03 再决定这些语义怎样由硬件路径实现。
\end{corebox}

忘掉公式后，尝试回答：
\begin{enumerate}
  \item 为什么指令长度、寄存器数量和立即数字段会竞争同一预算？
  \item 为什么 EA 不是 value？间接寻址为什么多一次访存？
  \item 为什么 32 位指令不能推出 32 位 PC 或物理地址？
  \item 为什么函数调用中的 caller-saved/callee-saved 属于 ABI 而非普适 ISA 自动动作？
  \item 给定 $x=a[i]+1$，能否从状态差独立重建 EA、load、加法和最终写回？
\end{enumerate}

\section{边界说明：哪些是稳定规则，哪些必须看题设}

本册中的状态差、有效地址、访问宽度和调用责任是跨架构都能使用的分析方法；MIPS 的固定格式、某套 ABI、某个处理器的节拍或寄存器分类只能作为题设例子。遇到题目没有说明的指令长度、PC 基准、异常返回点或字节序时，应明确写出假设，而不是从“32 位机器”等标签猜答案。

\section{补充细节：指令是有限编码上的状态转移契约}

指令题不是助记符翻译题，而是有限字段预算、地址形成和架构状态写集共同定义的状态转移。统一回接为
\[
\text{Instruction Format}
\to\text{Operand Interpretation}
\to\text{EA / Target Formation}
\to\text{Read/Write Set}
\to\text{Next Architectural State}.
\]

\subsection{一条指令必须携带的六类信息}

无论指令是定长还是变长，解码后都必须能回答六个问题：操作码是什么、源操作数在哪里、目的位置在哪里、操作数宽度是多少、立即数/位移如何解释、下一条 PC 如何产生。地址码个数（零地址、一地址、二地址、三地址）只是字段组织方式，不直接决定指令执行速度；同一语义可以通过更多临时寄存器或更复杂寻址表达。

机器字长、指令字长和存储字长也要分列。机器字长常影响寄存器/ALU 数据宽度，指令字长决定一次取指需要多少位，存储字长决定一个编址单元能放多少位；三者可以相同，也可以不同。地址线数决定可寻址单元数量，数据线数决定一次并行传输的位数，不能由“32 位机器”自动推出所有字段都是 32 位。

\begin{center}\small
\begin{tabularx}{\linewidth}{L{2.3cm}L{2.6cm}L{2.8cm}Y}
\toprule
字段/对象 & 它回答的问题 & 由什么约束 & 不能直接推出 \\
\midrule
opcode & 做什么 & 指令种类与控制动作 & 数据通路具体节拍 \\
register field & 对哪个架构寄存器做读写 & 寄存器数量的对数位宽 & 寄存器内部物理位置 \\
immediate/displacement & 常量或地址偏移如何编码 & 指令字长、符号扩展、范围 & 完整绝对地址 \\
addressing mode & 字段怎样解释为值/地址 & ISA 操作数语义 & Cache/TLB 命中结果 \\
width/sign & 读写多少位、怎样扩展 & 指令后缀或类型规则 & 总线一定只有同宽传输 \\
next-PC control & 顺序、分支或调用返回 & PC 基准与目标计算规则 & 预测器一定怎样实现 \\
\bottomrule
\end{tabularx}
\end{center}

\subsection{扩展操作码：剩余编码树而不是固定口诀}

设指令字长为 $L$，某一格式释放 $w$ 个地址字段位作为操作码，则本层最多有 $2^{L-w}$ 个前缀。若已经占用 $U$ 个普通编码，只有剩余的 $2^{L-w}-U$ 个前缀能够作为 escape；下一层若再释放 $q$ 位，就获得
\[
N_{next}=\bigl(2^{L-w}-U\bigr)2^q
\]
个候选编码。每层都要同时检查数量和前缀无歧义：一个 escape 不能既当普通指令又当下一层前缀，已经分配的前缀不能重复使用。

定长操作码易于并行译码和确定取指边界；变长/扩展操作码提高字段利用率和代码密度，却要求译码器先识别前缀，再知道后续字段的边界。题目改变地址字段数量、操作码层数或已用编码数时，必须重新画“剩余编码树”，不能套某道样题的答案。

\subsection{寻址方式：判据来自数据流，而不是名称}

有效地址统一写为
\[
EA=f(\text{base},\text{index},\text{scale},\text{disp},PC,M).
\]
立即数和寄存器寻址直接提供 value，因此通常没有数据访存；直接寻址让字段提供 EA；寄存器间接把寄存器内容当 EA；基址加偏移、变址和比例寻址把多个项组合成 EA；PC 相对寻址先确定 PC 基准再加符号扩展位移；间接寻址还需先从一个地址位置取出下一地址，再访问最终对象。

\begin{center}\small
\begin{tabularx}{\linewidth}{L{2.2cm}L{3.3cm}C{1.0cm}Y}
\toprule
方式 & 地址/值形成 & 数据访存（不含取指） & 解题检查 \\
\midrule
立即 & value = sext/zext(field) & 0 & 字段是值，不是地址 \\
寄存器 & $value=R_i$ & 0 & 不再对寄存器内容解引用 \\
直接 & EA = field & 1 & 字段单位、地址空间和宽度 \\
寄存器间接 & $EA=R_i$ & 1 & 指针值与被指对象分开 \\
基址/变址 & $EA=B+I\times K+d$ & 1 & $K$ 常为元素大小 \\
PC 相对 & $target=PC_{base}+sext(d)$ & 0 或取指路径 & PC 基准与位移单位 \\
一次间接 & $EA_1=M[field],\ value=M[EA_1]$ & 2 & 先取地址再取数据 \\
\bottomrule
\end{tabularx}
\end{center}

“地址的地址”是最容易被遗漏的层级：
\[
\text{value}=M[EA],\qquad
\text{pointer}=M[EA],\qquad
\text{next address}=M[M[EA]].
\]
访存次数必须从表达式层级数出来，不能从“间接/变址”四个字机械判断。

\subsection{有效地址计算：数组、相对分支和复合寻址}

对于元素大小为 $s$ 的一维数组，
\[
addr(a[i])=base(a)+i\times s.
\]
对 $N$ 列的行优先二维数组，
\[
addr(a[i][j])=base+(iN+j)\times s.
\]
如果字段是“元素下标”而非“字节偏移”，乘 $s$ 不能省略；如果 ISA 已经把位移定义为指令字或字对齐单位，则还要乘相应缩放因子。复合寻址先逐项写出 base、index、scale、displacement，再在最后一步得到 EA；不要把十六进制字段直接相加而跳过符号扩展。

PC 相对分支的安全表达是
\[
target=PC_{base}+\operatorname{sext}(disp)\times unit,
\]
其中 $PC_{base}$ 可能是当前指令地址、下一条指令地址或 ISA 规定的其他基准，$unit$ 可能是字节、字或指令槽。题目没有声明时，只能保留参数，不能从常见 MIPS/RISC-V 习惯强推。

\subsection{机器级表示：控制流和调用是架构状态图}

C 的选择、循环和过程调用都可还原成有限状态转移：比较产生条件，条件选择下一 PC，回边重复执行；调用把返回地址和参数交给被调用方，返回再把控制权和结果交回调用方。过程调用至少涉及：参数/返回值位置、返回地址保存、caller-saved 与 callee-saved 责任、栈帧布局、局部对象和对齐。

ISA 只提供跳转、寄存器读写和访存等原语；ABI 规定跨编译单元如何使用这些原语；编译器决定具体寄存器分配、是否内联、是否把变量放栈上。看到某个 x86/MIPS/RISC-V 寄存器名时，先判断它是题设架构事实还是 ABI 例子，不能把一套调用约定升级成通用 ISA。

\begin{methodbox}{递归/调用题的状态包}
对每次 call/return 明确写：入口 PC、返回 PC、参数位置、需要保存的寄存器、栈指针变化、返回值位置和退出后的下一 PC。这样可以把“画栈帧”转成可检查的 Read Set/Write Set，而不是背一张固定寄存器表。
\end{methodbox}

\subsection{汇编表示与目标文件：助记符不是最终机器状态}

汇编源到运行映像至少经历四个阶段：
\[
\text{Source}
\xrightarrow{compiler}\text{Assembly / Relocatable Object}
\xrightarrow{assembler}\text{Machine Object}
\xrightarrow{linker}\text{Executable Image}
\xrightarrow{loader}\text{Process Address Space}.
\]
编译器负责把高级语言语义映射为 ISA 可表达的指令和数据布局；汇编器解析助记符、伪指令、标签并生成机器码、符号表和重定位记录；链接器解析跨目标文件的符号引用、合并段并完成或保留重定位；装入器设置进程地址空间、入口 PC、栈和动态绑定。链接地址不是运行时物理地址，装入完成也不表示第一条指令已经提交。

x86 等变长指令的字节流必须先识别前缀、操作码、ModR/M、SIB、位移和立即数字段，再确定指令长度；Intel 与 AT\&T 语法的操作数顺序、立即数前缀和内存表达式书写不同，但最终机器语义应一致。伪指令可能展开成一条或多条真实指令，题目问“机器码长度”时必须确认是否已经展开。

\subsection{CISC 与 RISC：设计轴可以相关，但不能一一对应}

CISC 常允许较复杂的指令和寻址、较高代码密度与更复杂译码；RISC 倾向规整定长格式、较多寄存器、Load/Store 数据访问和硬件/编译器的清晰分工。它们描述的是一组设计取舍，不是控制器实现的逻辑定理：RISC 可以采用微程序控制，CISC 也可以在热点路径使用硬布线或微操作缓存。定长指令不等于指令格式种类少，Load/Store 也不等于程序没有内存访问，只是把访存集中到明确的 load/store 指令。

寄存器数量、寻址复杂度、代码密度、译码功耗、分支目标范围、编译器寄存器分配压力和流水线规整度应放在同一比较表中。任何“RISC 更快”或“CISC 指令更少所以总时间更短”的结论，都必须回到 $IC\times CPI\times T_{clk}$，并说明代码生成和实现模型。

\subsection{补充细节到 CO-03/CO-04 的交接}

\begin{center}\small
\begin{tabularx}{\linewidth}{L{2.7cm}L{3.0cm}L{3.0cm}Y}
\toprule
CO-02 输出 & 具体内容 & CO-03 消费方式 & CO-04 继续追踪 \\
\midrule
opcode/fields & 操作、寄存器、立即数、函数码 & 译码控制字和 Read/Write Set & 前端取指、译码时序 \\
EA/target & 地址生成式、PC 基准、访存次数 & ALU/MUX/寄存器微操作 & ready/need、分支预测与 flush \\
width/sign & 访问宽度、扩展、字节序、对齐 & load/store byte enable & 结构冲突、数据依赖 \\
call/return & 返回 PC、ABI 保存责任、栈帧 & 控制流与寄存器写回 & ordered commit、异常恢复 \\
relocation & 符号、重定位、装入地址 & 入口状态初始化 & 取指路径和 fault 重试 \\
\bottomrule
\end{tabularx}
\end{center}

\section{扩充后的陌生 ISA 复原检查}

面对一套没有学过的指令格式，仍按以下顺序复原：
\begin{enumerate}
  \item 先标总指令长度、编址单位、寄存器数量、字段边界和立即数符号规则；
  \item 对每条指令写 Read Set、Write Set、Next PC 和可能异常；
  \item 对每个内存操作写 EA 生成式、访问宽度、字节序和对齐条件；
  \item 若有扩展操作码，画剩余编码树并检查前缀唯一性；
  \item 若有函数调用，把 ISA 原语和 ABI 约定分成两层；
  \item 最后把状态包交给 CO-03 的 datapath/control 和 CO-04 的 timing/commit。
\end{enumerate}

\begin{methodbox}{本册扩充停止条件}
当读者能从任意新 ISA 的字段和题设行为重建“值/地址/地址的地址、访存次数、下一 PC、写回位置和异常边界”，CO-02 的机制层完成；具体架构的寄存器名称、汇编语法和编译器优化只作为题设实例，不再继续膨胀成跨架构结论。
\end{methodbox}

\end{document}
